%======================================================================
\section{Configuration lamps}\label{sec:group}
%======================================================================

Theorem~\ref{thm:criterion} turns $M$ copies of $S_3$ whose pairwise commutators have area at most $\Lambda$ into $\log\sigma_E(r)\ge M/16$ at
$r$ of order $M\Lambda$, so it approaches its ceiling $\exp(O(r))$ only when $\Lambda$ is tiny compared with $M$. The group of this section places
the copies of $S_3$ not on the points of Houghton's three rays but on their finitely supported binary configurations: Houghton's group moves the
points, finitely supported affine maps move the configurations, and one copy of $S_3$ sits on each configuration. Then $N$ points carry $2^N$
configurations, and hence $2^N$ copies of $S_3$, while every local move we use involves at most three points. Two of the $2^N$ copies are
compared by about $N^2$ local moves, each filled with the Dehn function of $H_3$ at scale $N$, so the area is polylogarithmic in the number of
copies.

Throughout, functions are composed from right to left, groups act on the left, $[x,y]=xyx^{-1}y^{-1}$, and $\log$ is to base two.

\paragraph{Houghton's group.} Let $Y=\{1,2,3\}\times\{1,2,\dots\}$; the point $(i,j)$ lies on ray $i$ at height $j$. Houghton's group
$H=H_3$ consists of the permutations $g$ of $Y$ that are eventually translations on each ray: there are integers $m_1(g),m_2(g),m_3(g)$ with
$g(i,j)=(i,j+m_i(g))$ for all large $j$. The map $g\mapsto(-m_2(g),-m_3(g))$ is a surjective homomorphism $H\to\Z^2$ whose kernel is the group
$\FSym(Y)$ of finitary permutations~\cite{Houghton1978}. Let $p$ be the permutation with $p(1,j)=(1,j+1)$, $p(2,1)=(1,1)$ and
$p(2,j)=(2,j-1)$ for $j\ge2$, fixing ray~$3$, and let $q$ be defined in the same way with rays $2$ and $3$ exchanged. They generate $H$
(Section~\ref{sec:houghton}), and their images in $\Z^2$ are $(1,0)$ and $(0,1)$.

\paragraph{The group.} Let $V=\F_2^{(Y)}$ be the space of finitely supported functions $Y\to\F_2$, the \emph{configurations}, with basis
$(e_y)_{y\in Y}$. The group $H$ acts on $V$ by permutation matrices, $g(e_y)=e_{g(y)}$, and we identify $H$ with its image in $\GL(V)$. Let
$\GL_{\rm fin}(Y)$ be the group of linear automorphisms of $V$ that fix all but finitely many basis vectors. It is the directed union of the finite
groups $\GL(\F_2^J)$, $J\subset Y$ finite, each extended by the identity on the other basis vectors; it is normalized by $H$; and
$H\cap\GL_{\rm fin}(Y)=\FSym(Y)$. Let $L=\GL_{\rm fin}(Y)\,H\le\GL(V)$. Write $\tau_u$ for the translation $v\mapsto v+u$ of $V$, let $A=V\rtimes L$
be the group of affine maps $v\mapsto g(v)+u$ with $g\in L$ and $u\in V$, and let
\[
\Gamma=\Bigl(\bigoplus_{v\in V}S_3\Bigr)\rtimes A,
\]
where $A$ permutes the summands through its action on $V$: writing $\lambda_v$ for the copy at $v$ of $\lambda\in S_3$, we have
$g\lambda_vg^{-1}=\lambda_{g(v)}$ for $g\in A$. Equivalently, $\Gamma$ is the group of permutations of $V\times S_3$ generated by the maps
$(v,s)\mapsto(g(v),s)$ for $g\in A$ and $(v,s)\mapsto(v,\lambda(v)s)$ for finitely supported $\lambda:V\to S_3$; this action is faithful. Groups obtained from a permutation action by adjoining the finitary linear or symmetric groups of the permuted
set, and lamps over it, appear as enrichments~\cite{Bradford2022,AlekseevBradford2026} and as halo products~\cite{GenevoisTessera2024}.

For distinct points $y,y'\in Y$ let $E_{y\leftarrow y'}$ be the transvection that sends $e_{y'}$ to $e_{y'}+e_y$ and fixes the other basis vectors.
For $g\in H$ and a linear map $M$ of $V$,
\begin{equation}\label{eq:affine-identities}
g\tau_{e_y}g^{-1}=\tau_{e_{g(y)}},\qquad gE_{y\leftarrow y'}g^{-1}=E_{g(y)\leftarrow g(y')},\qquad M\tau_uM^{-1}=\tau_{M(u)} .
\end{equation}
Linear maps fix $0$, so they commute with the copy of $S_3$ at $0$. On the first ray we abbreviate: for $x\ge1$ we write $x$ for the point $(1,x)$,
$e_x$ for $e_{(1,x)}$ and $\tau_x$ for $\tau_{e_x}$.

Let $a$ and $b$ be the transpositions $(1\,2)$ and $(2\,3)$ in the copy of $S_3$ at $0$, so that $ab$ has order three, and put
\[
t=\tau_1,\qquad d=E_{2\leftarrow1},\qquad c=d\,ab ;
\]
thus $t$ adds the point $1$ to a configuration, and $d$ adds the point $2$ to every configuration that contains the point $1$.

\begin{lemma}\label{lem:structure}
\begin{enumerate}[label=(\alph*),itemsep=1pt]
\item The element $c$ has order six, $c^3=d$ and $c^4=ab$; hence $b=ac^4$.
\item $\Gamma$ is generated by the five elements $p,q,t,c,a$.
\item Let $\pi:\Gamma\to\Z^2$ be the composite of $\Gamma\to A\to L\to L/\GL_{\rm fin}(Y)\cong H/\FSym(Y)\cong\Z^2$. Its kernel
$\bigl(\bigoplus_VS_3\bigr)\rtimes\bigl(V\rtimes\GL_{\rm fin}(Y)\bigr)$ is locally finite. Hence $\Gamma$ is locally finite-by-$\Z^2$, and in
particular elementary amenable.
\end{enumerate}
\end{lemma}

\begin{proof}
(a) The linear map $d$ fixes $0$, so it commutes with $ab$; since $d$ has order two and $ab$ order three, $c$ has order six, $c^3=d^3(ab)^3=d$ and
$c^4=d^4(ab)^4=ab$.

(b) Let $\Gamma'$ be the subgroup generated by the five elements. It contains $H$, because $H=\langle p,q\rangle$. Since $\FSym(Y)$ acts
$2$-transitively on $Y$, \eqref{eq:affine-identities} shows that $\Gamma'$ contains every $\tau_{e_y}$, hence $V$, and, since $d=c^3$, every
transvection $E_{y\leftarrow y'}$. Transvections generate every $\GL(\F_2^J)$, so $\GL_{\rm fin}(Y)\le\Gamma'$, and $A\le\Gamma'$. Finally
$b=ac^4\in\Gamma'$, and the copy of $S_3$ at $u\in V$ is $\tau_u\langle a,b\rangle\tau_u^{-1}$.

(c) Since $H\cap\GL_{\rm fin}(Y)=\FSym(Y)$, we have $L/\GL_{\rm fin}(Y)\cong H/\FSym(Y)$, and the kernel of $\pi$ is as stated. It is an
extension of $\bigoplus_VS_3$ by $V\rtimes\GL_{\rm fin}(Y)$, the directed union of the finite groups $\F_2^J\rtimes\GL(\F_2^J)$; both are locally
finite. An extension of a locally finite group $K$ by a locally finite group is locally finite: a finitely generated subgroup $\Lambda$ has finite
image modulo $K$, so $\Lambda\cap K$ has finite index in $\Lambda$, is finitely generated, and is finite. Locally finite groups are directed
unions of finite groups, and elementary amenable groups are closed under directed unions and extensions~\cite{Chou1980}.
\end{proof}

\subsection{A finite presentation}\label{sec:fp}

The certificate of Section~\ref{sec:certificate} is a finite set of relations and does not present $\Gamma$: the group it presents has
abelianization $\Z^2\times\Z/2\times\Z/4$, in which $c$ has order four (a computation with the exponent sums; for instance $a^2$,
$b^2=(ac^4)^2$ and $(ab)^3$ give $2a=8c=12c=0$), while in $\Gamma$ the element $c=d\,ab$ is a product of commutators, since
$d=[E_{2\leftarrow3},E_{3\leftarrow1}]$ and $ab$ is a three-cycle. The group itself is finitely presented.

\begin{proposition}[Finite presentation]\label{prop:fp}
The groups $L$, $A$ and $\Gamma$ are finitely presented.
\end{proposition}

\begin{proof}
\emph{Stabilizers.} For a finite set $D\subset Y$ the pointwise stabilizer $H_{(D)}$ of $D$ in $H$ is isomorphic to $H$: renumbering the points of each
ray outside $D$ in increasing order turns $Y\setminus D$ into three rays, and a permutation fixing $D$ is eventually a translation on the old rays
exactly when it is on the new ones. So $H_{(D)}$ is finitely generated. Since $H$ contains $\FSym(Y)$, it acts transitively on ordered $k$-tuples of
distinct points for every $k$.

\emph{A presentation lemma.} Let a finitely presented group $H$ act on a set $\mathcal X$ with finitely many orbits and finitely generated
stabilizers, and let $F(\mathcal X)$ be the free group on $\mathcal X$. Then $F(\mathcal X)\rtimes H$, with $H$ permuting the free generators, is
finitely presented: adjoin to a presentation of $H$ one representative $x_i$ of each orbit and the relations $[x_i,k]=1$ for $k$ in a finite
generating set of its stabilizer; then $hx_ih^{-1}$ depends only on $h\cdot x_i$, and the two groups are isomorphic (both are the iterated amalgam of $H$ with the groups
$\mathrm{Stab}(x_i)\times\langle x_i\rangle$ over the stabilizers $\mathrm{Stab}(x_i)$). Imposing, in addition, a set of
relations that is a finite union of $H$-orbits keeps the group finitely presented, since one representative of each orbit suffices.

\emph{A presentation of $A$.} For distinct $i,j\in Y$ let $x_{ij}$ be the transvection $e_j\mapsto e_j+e_i$, let $t_i=\tau_{e_i}$, and let
$\varsigma_{ij}\in\FSym(Y)\le H$ be the transposition of $i$ and $j$. Let $\widetilde A$ be the quotient of $F(\mathcal X)\rtimes H$, where
$\mathcal X=\{x_{ij}\}\sqcup\{t_i\}$ and $H$ relabels indices, by the relations
\begin{gather*}
x_{ij}^2=1,\qquad [x_{ij},x_{kl}]=1\ (i\ne l,\ j\ne k),\qquad [x_{ij},x_{jk}]=x_{ik}\ (i,j,k\text{ distinct}),\\
\varsigma_{ij}=x_{ij}x_{ji}x_{ij},\qquad t_i^2=1,\qquad [t_i,t_j]=1,\\
x_{ij}t_jx_{ij}^{-1}=t_it_j,\qquad x_{ij}t_kx_{ij}^{-1}=t_k\ (k\ne j).
\end{gather*}
They hold in $A$, so there is a surjection $\widetilde A\to A$. Each relation involves at most four points, so by transitivity the relations form
finitely many $H$-orbits; the stabilizers of $x_{ij}$ and $t_i$ are $H_{(\{i,j\})}$ and $H_{(\{i\})}$. Since $H$ is finitely presented~\eqref{eq:johnson}, $\widetilde A$ is finitely presented, and it
remains to show that $\widetilde A\to A$ is injective. The linear relations are Steinberg's~\cite{Milnor1971}; we give the short direct argument.

Fix a finite set $D\subset Y$, number its points $1,\dots,m$ (a numbering unrelated to heights), and let $U_D\le\widetilde A$ be generated by the $x_{ij}$ with $i<j$ in $D$. By the first three
relations, $x_{1m},\dots,x_{m-1,m}$ are commuting involutions, and they generate a subgroup normalized by the $x_{ij}$ with $i<j<m$; by induction
$|U_D|\le2^{m-1}2^{(m-1)(m-2)/2}=2^{m(m-1)/2}$, the order of the upper unitriangular group $\mathrm{UT}_m(\F_2)$, onto which $U_D$ maps. So
$U_D\cong\mathrm{UT}_m(\F_2)$. Let $W_D=\Sym(D)\le H$ and $\varsigma_i=\varsigma_{i,i+1}$. By the relation $\varsigma_{ij}=x_{ij}x_{ji}x_{ij}$ and relabeling, the
subgroup $G_D$ generated by all $x_{ij}$ with $i,j\in D$ contains $W_D$ and is generated by $U_D$ and the $\varsigma_i$. We claim $G_D=U_DW_DU_D$; since this
set contains $1$, it suffices that it is closed under right multiplication by $U_D$ and by each $\varsigma_i$. Every $b\in U_D$ is $x_{i,i+1}^\epsilon b_0$ with
$b_0$ in the subgroup generated by the other positive roots, the kernel of the entry $(i,i+1)$, which $\varsigma_i$ normalizes. So it suffices
to treat $wx_{i,i+1}\varsigma_i$ with $w\in W_D$; put $k=w(i)$ and $l=w(i+1)$. If $k<l$, then $wx_{i,i+1}\varsigma_i=x_{kl}w\varsigma_i\in U_DW_D$. If
$k>l$, then $w\varsigma_i=\varsigma_{kl}w$ and $x_{kl}\varsigma_{kl}=x_{lk}x_{kl}$, so $wx_{i,i+1}\varsigma_i=x_{kl}\varsigma_{kl}w=x_{lk}wx_{i,i+1}\in U_DW_DU_D$. The map $G_D\to\GL(\F_2^D)$ is onto, since transvections generate $\GL(\F_2^D)$. If $b_1wb_2$ maps to the identity of $\GL(\F_2^D)$, the permutation matrix of $w$
is upper unitriangular, so $w=1$ in $H$, and then $b_1b_2=1$ in $U_D$. Hence $G_D\cong\GL(\F_2^D)$. Every element of the subgroup $\widetilde J$
generated by all $x_{ij}$ lies in some $G_D$, so $\widetilde J\cong\GL_{\rm fin}(Y)$. By relabeling, every element of the subgroup generated by
$\widetilde J$ and $H$ has the form $jh$; if it maps to $1$, then $h\in\GL_{\rm fin}(Y)\cap H=\FSym(Y)$, so $h$ is a product of transpositions, each
equal in $\widetilde A$ to a word in the $x_{ij}$, and $jh\in\widetilde J$ maps to $1$, hence is $1$. Finally the $t_i$ generate a normal elementary
abelian subgroup, which maps isomorphically onto $V$ (it is a quotient of $\F_2^{(Y)}$, and $t_i\mapsto e_i$), and the linear part acts on it as on $V$, so every element of $\widetilde A$ is $t\ell$ with $t$ in that subgroup; if its image is
the identity affine map, evaluating at $0$ gives $t=1$, and then $\ell=1$. So $\widetilde A\cong A$. Omitting the $t_i$ gives $L$.

\emph{The lamps.} The group $A$ acts on $V$ doubly transitively: a translation carries $v$ to $0$, and a finitary linear map carries any nonzero
vector to $e_1$. The stabilizer of $0$ is $L$, which is generated by $p$, $q$ and $d$, because the $H$-conjugates of $d$ are all the transvections
$E_{y\leftarrow y'}$. So $A$ is finitely presented, its point stabilizers are finitely generated, and it has two orbits on $V\times V$. By de
Cornulier's criterion for permutational wreath products~\cite{Cornulier2006}, $\Gamma=S_3\wr_VA$ is finitely presented. Explicitly, if $\mathcal R_A$
presents $A$ on generators including $p,q,d,t$, then adding the generators $a,b$ and the relations $a^2$, $b^2$, $(ab)^3$, $[a,\ell]$ and $[b,\ell]$ for
$\ell\in\{p,q,d\}$, and $[u,tvt^{-1}]$ for $u,v\in\{a,b\}$, presents $\Gamma$: the lamp at $v\in V$ is well defined as a conjugate of
$\langle a,b\rangle$ because $L$ centralizes it, and double transitivity carries the commutation of the lamps at $0$ and $e_1$ to every pair.
\end{proof}

\subsection{An embedding in a Brin--Thompson group}\label{sec:3V}

Let $\mathfrak C=\{0,1\}^{\N}$ be Cantor space. For $n\ge1$, Brin's group $nV$~\cite{Brin2004} consists of the homeomorphisms $h$ of
$\mathfrak C^n$ given by a finite table of prefix replacements: there are two partitions of $\mathfrak C^n$ into the same number of rectangles
$u_1\mathfrak C\times\dots\times u_n\mathfrak C$, with $u_1,\dots,u_n$ finite binary words, matched in pairs, and $h$ maps each rectangle of the first
partition onto its partner by replacing the $n$ prefixes and keeping the $n$ suffixes. So $1V$ is Thompson's group $V$, and $nV$ embeds in $n'V$ for
$n\le n'$ by acting trivially on the additional coordinates. The groups $nV$ are finitely presented~\cite{HennigMatucci2012} and
simple~\cite{Brin2010}, and, as for $V$, it is not known whether they are sofic. (The letter $V$ in these names is not the space of configurations.)

\begin{proposition}\label{prop:3V}
$\Gamma$ embeds in $3V$, and hence in $nV$ for every $n\ge3$. Consequently $3V$ has a chunk $E_{3V}$ with
\[
\exp\Bigl(\frac{c\,r}{(\log r)^{12.75}}\Bigr)\ \le\ \sigma_{E_{3V}}(r)\ \le\ \exp(C\,r\log r)\qquad(r\ge2).
\]
\end{proposition}

\begin{proof}
\emph{Coding configurations.} For a finite set $S=\{j_1<\dots<j_m\}$ of positive integers, with $j_0=0$, let
\[
e(S)=\Bigl(\prod_{k=1}^m1\,0^{\,j_k-j_{k-1}-1}\,1\Bigr)\,0,
\]
so that $e(\emptyset)=0$, $e(\{1\})=110$ and $e(\{2\})=1010$. Read a word from the left: at a record boundary the letter $0$ ends the word and the
letter $1$ opens a record $10^k1$, which ends at its next letter $1$. So the words $e(S)$ form a prefix code. A finite word whose reading has not
ended can be extended to one that begins with some $e(S)$, by appending $0$ at a record boundary or $10$ inside a record. Hence the cylinders
$e(S)\mathfrak C$ are disjoint and their union is dense in $\mathfrak C$. For $v\in V$ and $i=1,2,3$ let $S_i(v)=\{j:v(i,j)=1\}$, let
\[
Z_v=e(S_1(v))\mathfrak C\times e(S_2(v))\mathfrak C\times e(S_3(v))\mathfrak C ,
\]
and let $\iota_v:\mathfrak C^3\to Z_v$ prefix the three codewords. The rectangles $Z_v$ are disjoint, and their union is dense in $\mathfrak C^3$.

\emph{Prepending a site.} Define $s_0,s_1:\mathfrak C\to\mathfrak C$ by $s_0(0\xi)=0\xi$, $s_0(1\xi)=10\xi$ and $s_1(\xi)=11\xi$. Their images
$0\mathfrak C\cup10\mathfrak C$ and $11\mathfrak C$ partition $\mathfrak C$, and, with $S+1=\{j+1:j\in S\}$,
\begin{equation}\label{eq:prepend}
s_0\bigl(e(S)\xi\bigr)=e(S+1)\,\xi,\qquad s_1\bigl(e(S)\xi\bigr)=e\bigl(\{1\}\cup(S+1)\bigr)\,\xi :
\end{equation}
an empty site in front lengthens the first gap by one, and an occupied site in front adds the record $11$.

\emph{The affine group.} For $b\in\{0,1\}$ put
\begin{gather*}
\hat P(x,s_b(y),z)=(s_b(x),y,z),\qquad \hat Q(x,y,s_b(z))=(s_b(x),y,z),\\
\hat T(s_b(x),y,z)=(s_{1-b}(x),y,z),
\end{gather*}
and let $\hat D$ fix $s_0(\mathfrak C)\times\mathfrak C^2$ pointwise and send $(11x,y,z)$ to $(11x',y,z)$, where $\hat T(x,y,z)=(x',y,z)$. These maps lie
in $3V$. With $\epsilon$ the empty word, $\hat P$ replaces the prefixes $(\epsilon,11,\epsilon)$, $(0,0,\epsilon)$, $(1,0,\epsilon)$, $(0,10,\epsilon)$ and
$(1,10,\epsilon)$ by $(11,\epsilon,\epsilon)$, $(0,0,\epsilon)$, $(10,0,\epsilon)$, $(0,1,\epsilon)$ and $(10,1,\epsilon)$; $\hat Q$ does the same with the roles
of the second and third coordinates exchanged; $\hat T$ exchanges the first-coordinate prefixes $0\leftrightarrow110$ and $10\leftrightarrow111$; and $\hat D$
exchanges $110\leftrightarrow11110$ and $1110\leftrightarrow11111$ and keeps $0$ and $10$. The permutation $p$ moves the first point of ray~$2$ to
the first point of ray~$1$ and shifts the other points of these rays, $q$ does the same with ray~$3$, $t$ changes the value of a configuration at
the point $1$, and $d$ adds the value at $1$ to the value at $2$. So~\eqref{eq:prepend} gives, for every $v\in V$,
\[
\hat P\iota_v=\iota_{p(v)},\qquad \hat Q\iota_v=\iota_{q(v)},\qquad \hat T\iota_v=\iota_{t(v)},\qquad \hat D\iota_v=\iota_{d(v)} .
\]
By the proof of Lemma~\ref{lem:structure}(b), $p,q,t,d$ generate $A$. If a word in them represents $g\in A$, the same word in $\hat P,\hat Q,\hat T,\hat D$ maps
$\iota_v(\xi)$ to $\iota_{g(v)}(\xi)$ for all $v$ and $\xi$; if $g=1$, it is the identity on the dense set $\bigcup_vZ_v$, hence the identity. So
$p,q,t,d\mapsto\hat P,\hat Q,\hat T,\hat D$ defines a homomorphism $\rho:A\to3V$ with
\begin{equation}\label{eq:3V-equivariant}
\rho(g)\,\iota_v=\iota_{g(v)}\qquad(g\in A,\ v\in V),
\end{equation}
and $\rho$ is injective: if $g\ne1$, then $g(v)\ne v$ for some $v$, and $\rho(g)$ maps $Z_v$ onto the disjoint rectangle $Z_{g(v)}$.

\emph{The lamps.} Let $\beta_1=0$, $\beta_2=10$ and $\beta_3=11$, and let $S_3$ act faithfully on $\mathfrak C^3$ by the prefix replacements
$\theta(\lambda)(\beta_i\xi,y,z)=(\beta_{\lambda(i)}\xi,y,z)$. For $\lambda\in S_3$ and $v\in V$ let $\rho(\lambda_v)$ be $\iota_v\theta(\lambda)\iota_v^{-1}$ on
$Z_v$ and the identity off $Z_v$; it lies in $3V$, since the complement of a rectangle is a finite union of rectangles. Copies at different
configurations have disjoint supports and commute, each copy is faithful, and by~\eqref{eq:3V-equivariant}
$\rho(g)\rho(\lambda_v)\rho(g)^{-1}=\rho(\lambda_{g(v)})$ for $g\in A$. These are the defining relations of $\Gamma$ as a semidirect product, so $\rho$
extends to a homomorphism $\rho:\Gamma\to3V$. Suppose that $\rho(\gamma)=1$, where $\gamma$ has lamp part $\lambda:V\to S_3$ and affine part $g$. The
lamps preserve every $Z_v$, so $\rho(g)$ does too, and $g(v)=v$ for all $v$; so $g=1$. Then $\rho(\gamma)$ restricts on $Z_v$ to
$\iota_v\theta(\lambda(v))\iota_v^{-1}$, so $\lambda(v)=1$ for all $v$. Hence $\rho$ is injective. Explicitly, $\rho(a)$ exchanges the rectangles
$00\mathfrak C\times0\mathfrak C\times0\mathfrak C$ and $010\mathfrak C\times0\mathfrak C\times0\mathfrak C$, and $\rho(c)=\hat D\rho(ab)$, where $\rho(ab)$
permutes the rectangles with prefixes $(00,0,0)$, $(010,0,0)$ and $(011,0,0)$ cyclically in this order and $\hat D$ is the identity on $Z_0$.

\emph{The chunk.} The sofic profile of a chunk depends only on its identity element and its partial multiplication
(Section~\ref{sec:criterion}), and an injective homomorphism preserves both. So the chunk $E_{3V}=\rho(E)$, for the chunk $E$ of
Theorem~\ref{thm:main-group}, has $\sigma_{E_{3V}}=\sigma_E$, and parts (a) and (b) of that theorem give the bounds.
\end{proof}

The script \texttt{verification/embedding3v\_check.py} composes the prefix tables exactly and checks the identities of the proof and
that the images of $p,q,t,c,a$ satisfy the $25$ relations of Table~\ref{tab:certificate}.

By Section~\ref{sec:thompson}, Thompson's group $V$ has chunks with finite superpolynomial profile, from its Houghton subgroups, and a chunk with
profile at least $\exp(c\,r/\log^2r)$, which is finite if $F$ is amenable. Proposition~\ref{prop:3V} gives finite, nearly exponential profile in
$3V$ with no assumption. We do not know whether $\Gamma$ embeds in $V$ or in $2V$ (Section~\ref{sec:open}).

%======================================================================
\section{Houghton inputs}\label{sec:houghton}
%======================================================================

We need three facts about $H$: short words for finite permutations, generators of point stabilizers with linear word length, and a bound on the
Dehn function. All transport below takes place among the first points of ray~$1$, so we only need words for permutations supported there. For
$n\ge1$ let $\Sigma_n\le H$ be the group of permutations of $Y$ supported in $\{(1,1),\dots,(1,n)\}$, identified with $\Sym(n)$; we write
$(x\ y)$ for the transposition of the points $x$ and $y$ of the first ray.

\paragraph{Conventions.} Put $\alpha=[p,q]$; a direct check shows $\alpha=(1\ 2)$. In~\cite{DouglasMghirbiA} permutations act on the right,
and Johnson's presentation of $H$ (see~\cite[Theorem~2.8]{Lee2012}) is written there in generators $a,b,\alpha$ that are the inverses, as maps,
of our $p,q,\alpha$. The map $g\mapsto g^{-1}$ is an isomorphism from the permutations of $Y$ with the product of~\cite{DouglasMghirbiA} to the
permutations of $Y$ under composition, and it sends $a,b,\alpha$ to $p,q,\alpha$. So a word in $a,b,\alpha$ represents $1$ in the conventions
of~\cite{DouglasMghirbiA} if and only if the same word in $p,q,\alpha$ represents $1$ here. Hence the relators of~\cite[Section~2]{DouglasMghirbiA},
read in our letters, give the presentation
\begin{equation}\label{eq:johnson}
P_J=\bigl\langle p,q,\alpha\bigm|\ \alpha^2,\ (\alpha p\alpha p^{-1})^3,\ [\alpha,p^2\alpha p^{-2}],\ [p,q]\alpha^{-1},\ p\alpha p^{-1}q\alpha^{-1}q^{-1}\bigr\rangle
\end{equation}
of $H$; in particular $H=\langle p,q\rangle$. Renaming the letters is an isomorphism of free groups that carries the relators
of~\cite{DouglasMghirbiA} to those of $P_J$, so every area computed in~\cite{DouglasMghirbiA} for Johnson's presentation holds verbatim for $P_J$.

Eliminating $\alpha$ by the relator $[p,q]\alpha^{-1}$, a Tietze transformation, gives the presentation
\begin{equation}\label{eq:PH}
P_H=\bigl\langle p,q\bigm|\ s^2,\ (sps\,p^{-1})^3,\ [s,p^2sp^{-2}],\ psp^{-1}qs^{-1}q^{-1}\bigr\rangle,\qquad s=[p,q],
\end{equation}
of $H$. Let $\delta_H$ be its Dehn function: $\delta_H(n)$ is the largest area, with respect to the four relators of $P_H$, of a word of length at most
$n$ in $p^{\pm1},q^{\pm1}$ that represents $1$ in $H$. It is nondecreasing.

\begin{lemma}[Finite permutations]\label{lem:transpositions}
Put $\sigma_k=p^k\alpha p^{-k}$ for $k\ge0$.
\begin{enumerate}[label=(\alph*),itemsep=1pt]
\item $\sigma_k=(k+1\ \,k+2)$, and for $1\le x<y$ the word
$\omega_{x,y}=p^{x-1}(\alpha p)^{y-x-1}(\alpha p^{-1})^{y-x-1}\alpha\,p^{1-x}$ represents $(x\ y)$. With $\alpha=pqp^{-1}q^{-1}$ it has length at most
$10y-8x-8<10y$ in $p,q$.
\item Every element of $\Sigma_n$ that moves $m\ge2$ points is represented by a word in $p,q$ of length less than $10(m-1)n$.
\item Let $x_1,\dots,x_l\in\{1,\dots,n\}$ be distinct. Put $\pi_0=1$ and, for $m=1,\dots,l$, $w_m=\pi_{m-1}^{-1}(x_m)$ and
$\pi_m=\pi_{m-1}(m\ w_m)$, where $(m\ m)=1$. Then $\pi[x_1,\dots,x_l]:=\pi_l$ sends $m$ to $x_m$ for $m\le l$, it is a product of at most $l$
transpositions of points in $\{1,\dots,n\}$, and the product of the corresponding words $\omega$ represents it with length less than $10ln$.
\end{enumerate}
\end{lemma}

\begin{proof}
(a) Since $p^k$ maps $(1,1),(1,2)$ to $(1,k+1),(1,k+2)$, $\sigma_k$ is the stated transposition. Hence
$(x\ y)=\sigma_{x-1}\sigma_x\cdots\sigma_{y-3}\,\sigma_{y-2}\,\sigma_{y-3}\cdots\sigma_{x-1}$, by induction on $y-x$ from
$(x\ y)=\sigma_{x-1}(x{+}1\ \,y)\sigma_{x-1}$. In the product the powers of $p$ telescope: $\sigma_k\sigma_{k+1}=p^k\alpha p\alpha p^{-k-1}$ and
$\sigma_{k+1}\sigma_k=p^{k+1}\alpha p^{-1}\alpha p^{-k}$, which gives $\omega_{x,y}$ after free reduction. Its length is
$2(x-1)+10(y-x-1)+4$.

(b) A permutation moving $m$ points is a product of at most $m-1$ transpositions of those points, each of length less than $10n$ by (a).

(c) Since $\pi_{m-1}(l)=x_l\ne x_m$ for $l<m$, we have $w_m\ge m$, so $(m\ w_m)$ fixes $1,\dots,m-1$ and $\pi_m(l)=x_l$ for $l\le m$. Each factor
has length less than $10n$ by~(a).
\end{proof}

\begin{lemma}[Point stabilizers]\label{lem:stabilizers}
For $k\ge1$ let $H^{(k)}$ be the pointwise stabilizer in $H$ of $(1,1),\dots,(1,k)$, put $\zeta_k=\sigma_{k-1}\cdots\sigma_1\sigma_0$ and
$K_k=(\zeta_kp,\ \zeta_kq)$.
\begin{enumerate}[label=(\alph*),itemsep=1pt]
\item There is an isomorphism $H\to H^{(k)}$ that sends $p,q,\alpha$ to $\zeta_kp,\zeta_kq,\sigma_k$ and the transposition $(x\ y)$ to
$(x{+}k\ \ y{+}k)$. In particular $\zeta_kp$ and $\zeta_kq$ generate $H^{(k)}$.
\item Every element of $\Sigma_n\cap H^{(k)}$ that moves $m\ge2$ points is represented by a word of length less than $10(m-1)n$ in $\zeta_kp$ and
$\zeta_kq$.
\end{enumerate}
In $p,q$ the words $\zeta_1p=[p,q]\,p$ and $\zeta_2p=p[p,q]p^{-1}[p,q]\,p$ have length at most $5$ and $11$, and likewise for $q$.
\end{lemma}

\begin{proof}
(a) Let $\theta_k:Y\to Y\setminus\{(1,1),\dots,(1,k)\}$ send $(1,j)$ to $(1,j+k)$ and fix rays $2$ and $3$. For $g\in H$ let $g^{[k]}$ act as
$\theta_kg\theta_k^{-1}$ on the image of $\theta_k$ and fix $(1,1),\dots,(1,k)$. Then $g^{[k]}$ is eventually a translation on each ray, and
$g\mapsto g^{[k]}$ is an isomorphism $H\to H^{(k)}$, since an element of $H^{(k)}$ restricts to a permutation of the image of $\theta_k$. Clearly
$(x\ y)^{[k]}=(x{+}k\ \ y{+}k)$, so $\alpha^{[k]}=\sigma_k$. The element $\zeta_k$ is the cycle $(1,1)\mapsto(1,k+1)\mapsto(1,k)\mapsto\dots\mapsto(1,2)\mapsto(1,1)$.
Hence $\zeta_kp$ fixes $(1,j)$ for $j\le k$, maps $(1,j)$ to $(1,j+1)$ for $j>k$, maps $(2,1)$ to $(1,k+1)$, and agrees with $p$ elsewhere; that is,
$\zeta_kp=p^{[k]}$. Likewise $\zeta_kq=q^{[k]}$.

(b) An element $g\in\Sigma_n\cap H^{(k)}$ moving $m$ points is $g'^{[k]}$ for some $g'\in\Sigma_{n-k}$ moving $m$ points. Write $g'$ as in
Lemma~\ref{lem:transpositions}(b) and substitute $\zeta_kp,\zeta_kq$ for $p,q$.
\end{proof}

\begin{proposition}[Dehn function]\label{prop:dehn}
Let $\delta_J$ be the Dehn function of $P_J$. Then $\delta_H\le\delta_J$, and:
\begin{enumerate}[label=(\alph*),itemsep=1pt]
\item there is an absolute constant $K$ with $\delta_H(n)\le e^{Kn}$ for $n\ge1$;
\item there is an absolute constant $C$ with $\delta_H(n)\le Cn^{\,p_H}$ for $n\ge1$, where $p_H=\log107+4\approx10.74$.
\end{enumerate}
\end{proposition}

\begin{proof}
Let $w$ be a word in $p,q$ of length at most $n$ representing $1$. In the free group on $p,q,\alpha$ it is a product of at most $\delta_J(n)$
conjugates of relators of $P_J$ and their inverses. Apply the homomorphism to the free group on $p,q$ that fixes $p,q$ and sends $\alpha$ to
$[p,q]$. It fixes $w$, sends $[p,q]\alpha^{-1}$ to the trivial word, and sends each of the other four relators to a relator of $P_H$. Hence
$\Area_{P_H}(w)\le\delta_J(n)$.

(a) Lee proved that $H_3$ satisfies an exponential isoperimetric inequality~\cite{Lee2012}. The Dehn functions $\delta,\delta'$ of two finite
presentations of the same group are equivalent: $\delta(n)\le C'\delta'(C'n)+C'n+C'$ for a constant $C'$~\cite{Bridson2002}. This bound
preserves the class of functions bounded by $e^{Kn}$.

(b) By~\cite[Theorem~3.3]{DouglasMghirbiA}, $\delta_J(n)\le400n^4(1+A^*(4n+2))$, where the far-commutator area satisfies
$A^*(M)\le135M^{\log107}$~\cite[Theorem~3.2]{DouglasMghirbiA}; so $\delta_J(n)=O(n^{\log107+4})$. By the discussion of conventions above this
holds for $P_J$ in our letters, and $\delta_H\le\delta_J$.
\end{proof}

Part~(a) uses only the published exponential bound, and part~(b) our polynomial one. Either enters the lower bounds below only through
$\delta_H$ at scale $N$, the number of points that carry the configurations.

%======================================================================
\section{A certificate of 25 relations and local fillings}\label{sec:certificate}
%======================================================================

All fillings in $\Gamma$ are built from one finite set $R$ of words in the five generators of Lemma~\ref{lem:structure}, each representing $1$
in $\Gamma$. The set $R$ does not present $\Gamma$ (Section~\ref{sec:fp}); areas are always taken with respect to $R$.

\paragraph{Abbreviations.} Put
\[
s=[p,q],\qquad B=psp^{-1},\qquad t_m=p^{m-1}\,t\,p^{1-m}\quad(m=1,2,3),\qquad d=c^3,\qquad b=ac^4 .
\]
By Lemma~\ref{lem:transpositions}, $s$ and $B$ are the transpositions $(1\ 2)$ and $(2\ 3)$, and $t_m$ represents $\tau_m$, since
$p^{m-1}(1,1)=(1,m)$. By Lemma~\ref{lem:stabilizers}, with $\zeta_1=s$ and $\zeta_2=Bs$, the pairs
\[
K_1=(sp,\ sq),\qquad K_2=(Bsp,\ Bsq)
\]
generate the pointwise stabilizers $H_1$ of the point $1$ and $H_2$ of the points $1,2$. The elements $t$ and $d$ have the \emph{prototypes}
$P_t=\{1\}$ and $P_d=\{1,2\}$, the points on which they depend: by~\eqref{eq:affine-identities}, an element of $H$ that fixes $P_f$ pointwise
commutes with $f$.

\begin{table}[ht]
\centering\small
\begin{tabular}{@{}l p{0.64\textwidth} r r@{}}
\toprule
 & relations & number & length\\
\midrule
\textsf{H} & $s^2$,\ $(sps\,p^{-1})^3$,\ $[s,p^2sp^{-2}]$,\ $psp^{-1}qs^{-1}q^{-1}$ & 4 & 64\\
\textsf{Q} & $t^2$ & 1 & 2\\
\textsf{S} & $[t,k]$ for $k\in K_1$;\ \ $[d,k]$ for $k\in K_2$ & 4 & 56\\
\textsf{Z} & $[g,x]$ for $g\in\{p,q\}$, $x\in\{a,b\}$;\ \ $[d,a]$ & 5 & 40\\
\textsf{C} & $[t_1,t_2]$ & 1 & 8\\
\textsf{A} & $d\,t_1d^{-1}(t_1t_2)^{-1}$,\ $[d,t_2]$,\ $[d,t_3]$ & 3 & 39\\
\textsf{L} & $a^2$,\ $b^2$,\ $(ab)^3$ & 3 & 30\\
\textsf{P} & $[x,tyt^{-1}]$ for $x,y\in\{a,b\}$ & 4 & 64\\
\midrule
 & total & 25 & 303\\
\bottomrule
\end{tabular}
\caption{The certificate $R$: eight families of relations in the five generators $p,q,t,c,a$, written with the abbreviations above. Each
relation is taken to be the freely reduced word obtained by expanding the abbreviations; lengths are those of these words, and the longest has
length $24$.}
\label{tab:certificate}
\end{table}

\begin{lemma}\label{lem:validity}
Every word of Table~\ref{tab:certificate} represents~$1$ in $\Gamma$.
\end{lemma}

\begin{proof}
\textsf{H}: the words of~\eqref{eq:PH} hold in $H$.
\textsf{Q}: $V$ has exponent two.
\textsf{S}: the elements of $H_1$ fix the point $1$, so they commute with $t=\tau_1$; those of $H_2$ fix the points $1$ and $2$, so they commute
with $d=E_{2\leftarrow1}$ (Lemma~\ref{lem:structure}(a) and~\eqref{eq:affine-identities}).
\textsf{Z}: $p$, $q$ and $d$ are linear, so they commute with $a$ and with $b=ac^4=a\cdot ab$, which lie in the copy of $S_3$ at~$0$.
\textsf{C}: translations commute.
\textsf{A}: $dt_md^{-1}=\tau_{d(e_m)}$, and $d(e_1)=e_1+e_2$, while $d$ fixes $e_2$ and $e_3$.
\textsf{L}: these are relations of $S_3$.
\textsf{P}: $tyt^{-1}$ lies in the copy of $S_3$ at $e_1$, which commutes with the copy at~$0$.
\end{proof}

The $25$ freely reduced words are listed in the file \texttt{verification/relators1d.txt}. The script \texttt{verification/relators1d\_check.py}
rebuilds each word from the abbreviations and checks it, independently of Lemma~\ref{lem:validity}, in the faithful action of $\Gamma$ on
$V\times S_3$. For a word $w$ of length $l$ it checks that the exponent sums of $p$ and $q$ vanish; that $w$ fixes the empty configuration and
every configuration $\{y\}$ with $y$ of height at most $l+2$ or far out on a ray; and that it fixes $(v,s)$ for every $s\in S_3$ at every
configuration $v$ at which one of its lamp letters acts. These checks are exhaustive: while $w$ is applied, a point of height greater than $l+2$
never reaches the points $1$ and $2$ and is moved only by translations along its ray, so the checks determine the affine part of $w$, and then
the lamps of $w$ sit at the configurations just listed. The script also checks random configurations, that the abbreviations act as stated, and,
as a negative control, that deleting any single letter from a sample of eleven relations produces a word that acts nontrivially.

\subsection{Local fillings}

Let $n\ge3$. We use three elementary facts about areas. The area is invariant under free equality, inversion and conjugation, and subadditive
under products. If $N$ represents~$1$, then $\mu N\nu=(\mu N\mu^{-1})\mu\nu$ for all words $\mu,\nu$; so replacing a subword $U$ by $U'$ changes
the area by at most $\Area_R(UU'^{-1})$. And $[w,z_1z_2]=[w,z_1]\cdot z_1[w,z_2]z_1^{-1}$ and $[w,z^{-1}]=z^{-1}[w,z]^{-1}z$, so
$\Area_R([w,z])$ is at most the sum of the areas $\Area_R([w,z_l])$ over the letters $z_l$ of any factorization $z=z_1^{\pm1}\cdots z_m^{\pm1}$.

\paragraph{Canonical transporters.} For points $i\ne j$ in $\{1,\dots,n\}$ put
\[
U_i=\pi[i],\qquad \tau_i=U_i\,t\,U_i^{-1},\qquad U_{ij}=\pi[j,i],\qquad X_{ij}=U_{ij}\,d\,U_{ij}^{-1},
\]
with the words of Lemma~\ref{lem:transpositions}(c). We use $\tau_i$ both for this word and for the translation by $e_i$ that it represents, since
$U_i(1)=i$; and $X_{ij}$ represents $E_{i\leftarrow j}$, since $U_{ij}$ maps $1$ to $j$ and $2$ to $i$. By Lemma~\ref{lem:transpositions},
$|U_i|<10n$ and $|U_{ij}|<20n$, and $U_i$ ($U_{ij}$) is a product of at most one (two) transpositions.

The next lemma says that it does not matter which short word transports a prototype: two such words differ by an element of a stabilizer,
which commutes with the transported element by the family~\textsf{S}. We allow the transporting word to end with a power of $p$, because the
relations of Table~\ref{tab:certificate} use $t_2=ptp^{-1}$ and $t_3=p^2tp^{-2}$.

\begin{lemma}[Change of transporter]\label{lem:exchange}
Let $n\ge3$ and $f\in\{t,d\}$. Let $U$ and $W_0$ be words in $p^{\pm1},q^{\pm1}$ of length less than $30n$ that represent elements of $\Sigma_n$,
such that $W_0^{-1}U$ moves at most $10$ points, and let $W=W_0p^j$, where $j=0$ if $f=d$ and $j\in\{0,1,2\}$ if $f=t$. If $W$ and $U$ agree
on $P_f$, then
\[
\Area_R\bigl(WfW^{-1}\cdot Uf^{-1}U^{-1}\bigr)\ \le\ \mathcal E(n):=2\,\delta_H(1500n)+130n .
\]
\end{lemma}

\begin{proof}
Let $k=|P_f|$ and let $g$ be the element represented by $W^{-1}U$; it fixes $P_f$ pointwise, so $g\in H^{(k)}$. If $j=0$, then $g\in\Sigma_n$ moves at
most $10$ points, and by Lemma~\ref{lem:stabilizers}(b) it is represented by a word $Z$ with fewer than $90n$ letters from $K_k^{\pm1}$. If
$j\ge1$, then $f=t$, $k=1$, and $g=(sp)^{-j}g'$ with $g'=(sp)^jp^{-j}\cdot W_0^{-1}U$. Here $(sp)p^{-1}=s$ and $(sp)^2p^{-2}=sB$ lie in
$\Sigma_3$, so $g'\in\Sigma_n$ moves at most $13$ points; and $g'$ fixes the point $1$, since $g$ and $sp$ do. By Lemma~\ref{lem:stabilizers}(b),
$g'$ is represented by a word $Z'$ with fewer than $120n$ letters from $K_1^{\pm1}$, and we put $Z=(sp)^{-j}Z'$. In both cases $Z$ has fewer than
$122n$ letters from $K_k^{\pm1}$, each of length at most $11$ in $p,q$.

Let $N=W^{-1}UZ^{-1}$. It represents $1$ in $H$ and has length less than $30n+2+30n+1342n\le1500n$, so by the family \textsf{H} its area is at
most $\delta_H(1500n)$. By the family \textsf{S}, $\Area_R([f,Z])<122n$. Since $U=WNZ$ freely,
\[
WfW^{-1}\cdot Uf^{-1}U^{-1}=W\,(fNf^{-1})\,[f,Z]\,N^{-1}\,W^{-1},
\]
whose area is at most $2\delta_H(1500n)+122n\le\mathcal E(n)$.
\end{proof}

\begin{lemma}[Local fillings]\label{lem:local-filling}
Let $n\ge3$ and
\[
L(n)=10\,\delta_H(1500n)+700\,n .
\]
Let $i,j,k\in\{1,\dots,n\}$ be points with $i\ne j$, and $x,y\in\{a,b\}$, where $b$ stands for the word $ac^4$. The following words represent $1$,
and each has area at most $L(n)$ with respect to $R$:
\begin{enumerate}[label=(\alph*),itemsep=1pt]
\item $\tau_i^2$, of area $1$, and $[\tau_i,\tau_j]$;
\item $X_{ij}\tau_kX_{ij}^{-1}\tau_k^{-1}$ if $k\ne j$, and $X_{ij}\tau_jX_{ij}^{-1}\tau_i^{-1}\tau_j^{-1}$;
\item $[X_{ij},x]$, of area less than $40n+1$;
\item $[x,\tau_iy\tau_i^{-1}]$, of area less than $40n+1$.
\end{enumerate}
\end{lemma}

\begin{proof}
All the words represent $1$ by~\eqref{eq:affine-identities}. Call replacing a subword $UfU^{-1}$ by $WfW^{-1}$ as in Lemma~\ref{lem:exchange} a
\emph{change of transporter}; it costs at most $\mathcal E(n)$, and $5\mathcal E(n)+1\le L(n)$. Below, $W_0$ is always a product of at most three
transpositions, of length less than $30n$, and $U\in\{U_i,U_j,U_k,U_{ij}\}$ is a product of at most two, so $W_0^{-1}U$ moves at most $10$ points.

(a) $\tau_i^2=U_it^2U_i^{-1}$ freely, a conjugate of the relation \textsf{Q}. For $[\tau_i,\tau_j]$ let $W_0=\pi[i,j]$. Replace both occurrences
of $\tau_i^{\pm1}$ by $(W_0tW_0^{-1})^{\pm1}$, and both occurrences of $\tau_j^{\pm1}$ by $(W_0t_2W_0^{-1})^{\pm1}=(W_0p\,t\,p^{-1}W_0^{-1})^{\pm1}$.
These are four changes of transporter, since $W_0$ and $U_i$ both map $1$ to $i$, and $W_0p$ and $U_j$ both map $1$ to $j$. The result
$W_0[t_1,t_2]W_0^{-1}$ is a conjugate of the relation \textsf{C}, so the area is at most $4\mathcal E(n)+1$.

(b) Let $W_0=\pi[j,i,k]$, or $\pi[j,i]$ if $k\in\{i,j\}$, so that $W_0$ maps $1$ to $j$, $2$ to $i$, and $3$ to $k$ when $k\notin\{i,j\}$; let
$m\in\{1,2,3\}$ be the point with $W_0(m)=k$, so that $m=1$ exactly when $k=j$. Replace $X_{ij}^{\pm1}$ by $(W_0dW_0^{-1})^{\pm1}$, a change of
transporter because $W_0$ and $U_{ij}$ agree on $P_d=\{1,2\}$. Replace $\tau_k^{\pm1}$ by $(W_0t_mW_0^{-1})^{\pm1}=(W_0p^{m-1}\,t\,p^{1-m}W_0^{-1})^{\pm1}$,
and $\tau_i^{\pm1}$ by $(W_0t_2W_0^{-1})^{\pm1}$, again changes of transporter. If $k\ne j$, four of them turn the word into
$W_0(dt_md^{-1}t_m^{-1})W_0^{-1}$ with $m\in\{2,3\}$, a conjugate of a relation \textsf{A}. If $k=j$, five of them turn it into
$W_0(dt_1d^{-1}t_2^{-1}t_1^{-1})W_0^{-1}$, again a conjugate of a relation \textsf{A}. So the area is at most $5\mathcal E(n)+1$.

(c) Moving $x$ across a letter of $U_{ij}$ costs one relation \textsf{Z}, and moving it across $d=c^3$ costs one relation $[d,a]$, since
$[c^3,ac^4]=[c^3,a]$ freely. Hence $\Area_R([X_{ij},x])\le2|U_{ij}|+1<40n+1$.

(d) Freely the word is $U_i[U_i^{-1}xU_i,\,t(U_i^{-1}yU_i)t^{-1}]U_i^{-1}$. Replacing the two occurrences of $U_i^{-1}x^{\pm1}U_i$ and the two of
$U_i^{-1}y^{\pm1}U_i$ by $x^{\pm1}$ and $y^{\pm1}$ costs at most $4|U_i|$ relations \textsf{Z} and leaves $U_i[x,tyt^{-1}]U_i^{-1}$, a conjugate of a
relation \textsf{P}; so the area is at most $4|U_i|+1<40n+1$.
\end{proof}

%======================================================================
\section{Pair areas and the profile}\label{sec:profile}
%======================================================================

We now bring together two of the $2^N$ copies of $S_3$ supported on the first $N$ points of ray~$1$. Their transport words are products of the
$\tau_i$; the commutator of two copies reduces to the copy at a single point after clearing the difference of the two configurations by at most
$N$ transvections. Clearing naively would let the translation word grow; we keep it normalized after every step, so that it never has more than
$N$ factors, and each step costs $O(N)$ local fillings.

\begin{theorem}[Pair areas]\label{thm:pair-area}
Let $N\ge1$ and $n=\max(3,N)$. For a configuration $u$ supported in the first $N$ points of ray~$1$ let $T_u=\tau_{i_1}\cdots\tau_{i_m}$,
where $i_1<\dots<i_m$ are the points of the support of $u$. Then $T_u$ represents $\tau_u$, so $T_uaT_u^{-1}$ and $T_ubT_u^{-1}$ generate the copy
of $S_3$ at $u$; $|T_u|\le21Nn$; and for distinct such $u,w$ and $x,y\in\{a,b\}$, with $b$ the word $ac^4$,
\[
\Area_R\bigl([T_uxT_u^{-1},\,T_wyT_w^{-1}]\bigr)\ \le\ \Lambda_N:=20N^2L(n)\ \le\ C_0N^2\bigl(\delta_H(C_0N)+N\bigr)
\]
with $C_0=4.5\cdot10^4$. Consequently there are absolute constants $K'$ and $C'$ with
\[
\Lambda_N\le e^{K'N}\qquad\text{and}\qquad \Lambda_N\le C'N^{2+p_H}\qquad(N\ge1),\qquad p_H=\log107+4 .
\]
\end{theorem}

\begin{proof}
Translations commute, so $T_u$ represents $\tau_u$, and $|T_u|\le N(2\cdot10n+1)\le21Nn$. Throughout, $L=L(n)$; by
Lemma~\ref{lem:local-filling} the local identities (a)--(d) have area at most $L$.

\emph{Step 1: one translation word.} Freely, $[T_uxT_u^{-1},T_wyT_w^{-1}]=T_u[x,ZyZ^{-1}]T_u^{-1}$ with $Z=T_u^{-1}T_w$. Replace each of the at
most $N$ factors $\tau_i^{-1}$ of $T_u^{-1}$ by $\tau_i$ (one relation \textsf{Q} each), sort the resulting product of at most $2N$ factors
$\tau_i$ into increasing order by at most $\binom{2N}2$ swaps of adjacent distinct factors (Lemma~\ref{lem:local-filling}(a)), and cancel the at
most $N$ adjacent equal pairs (one relation \textsf{Q} each). The result is $T_{u_0}$ with $u_0=u+w\ne0$, and
$\Area_R(ZT_{u_0}^{-1})\le2N+N(2N-1)L\le2N^2L$. Since $Z^{\pm1}$ occurs four times in $[x,ZyZ^{-1}]$,
\[
\Area_R\bigl([T_uxT_u^{-1},T_wyT_w^{-1}]\bigr)\le8N^2L+\Area_R\bigl([x,T_{u_0}yT_{u_0}^{-1}]\bigr).
\]

\emph{Step 2: normalized clearing.} Let $s\le N$ be the size of the support of $u_0$, choose a point $j$ in it, and list the other points of the
support as $i_1,\dots,i_{s-1}$. Put $G_l=X_{i_lj}\cdots X_{i_1j}$ and $u_l=u_0+e_{i_1}+\dots+e_{i_l}$, whose support is that of $u_0$ with
$i_1,\dots,i_l$ removed; it still contains $j$, and $u_{s-1}=e_j$. We claim
\[
\Area_R\bigl(G_lT_{u_0}G_l^{-1}\,T_{u_l}^{-1}\bigr)\le2lsL\qquad(0\le l\le s-1).
\]
For $l=0$ there is nothing to prove. Given the claim for $l$, write $X=X_{i_{l+1}j}$ and $i=i_{l+1}$. Freely
$G_{l+1}T_{u_0}G_{l+1}^{-1}=X(G_lT_{u_0}G_l^{-1})X^{-1}$, which differs from $XT_{u_l}X^{-1}$ by a conjugate of a word of area at most $2lsL$. Next,
$XT_{u_l}X^{-1}$ is freely the product of the words $X\tau_{i'}X^{-1}$ over the support of $u_l$, in order. By Lemma~\ref{lem:local-filling}(b),
at a cost of at most $sL$, each of these becomes $\tau_{i'}$, except that $X\tau_jX^{-1}$ becomes $\tau_j\tau_i$. The new factor $\tau_i$ has a
partner $\tau_i$ in $T_{u_l}$, because $i$ lies in the support of $u_l$; at most $s-1$ swaps (Lemma~\ref{lem:local-filling}(a)) make the two
adjacent, and one relation \textsf{Q} cancels them. What remains is $T_{u_{l+1}}$. This step costs at most $sL+(s-1)L+1\le2sL$, which proves the
claim. The translation word is renormalized after every transvection and never has more than $s\le N$ factors.

\emph{Step 3: a single point.} Let $G=G_{s-1}$ and let $M=GT_{u_0}G^{-1}\tau_j^{-1}$, which has area at most $2s(s-1)L\le2N^2L$ by Step~2. Freely
$T_{u_0}=G^{-1}M\tau_jG$. Substituting this into the four occurrences of $T_{u_0}^{\pm1}$ in $[x,T_{u_0}yT_{u_0}^{-1}]$ and deleting the four
occurrences of $M^{\pm1}$ costs at most $8N^2L$. By Lemma~\ref{lem:local-filling}(c), $\Area_R([G,z])\le(s-1)(40n+1)$ for $z\in\{a,b\}$, and the
same bound holds for $[G,z^{-1}]$ and $[z,G^{-1}]$. Moving $y^{\pm1}$ across $G$ and $G^{-1}$, and then $x^{\pm1}$, leaves
$G^{-1}[x,\tau_jy\tau_j^{-1}]G$, at a cost of at most $4(s-1)(40n+1)$; and $[x,\tau_jy\tau_j^{-1}]$ has area at most $40n+1$ by
Lemma~\ref{lem:local-filling}(d). Altogether, since $L\ge700n$,
\[
\Area_R\bigl([x,T_{u_0}yT_{u_0}^{-1}]\bigr)\le8N^2L+4N(40n+1)+40n+1\le9N^2L,
\]
and with Step~1 the area is at most $17N^2L\le\Lambda_N$.

\emph{Growth.} Since $n\le3N$, $L(n)\le10\delta_H(4500N)+2100N$, which gives the first bound on $\Lambda_N$. With Proposition~\ref{prop:dehn}(a)
it is at most $e^{K'N}$ for a suitable $K'$. With Proposition~\ref{prop:dehn}(b) it is at most $C_0N^2(C(C_0N)^{p_H}+N)\le C'N^{2+p_H}$, since $p_H>1$.
\end{proof}

The factor $N^2$ counts local moves: at most $N$ transvections, each followed by at most $N$ commutations. The Dehn function of $H$ enters only
through the local fillings, at the scale $N$ of the transport.

\subsection{Affine hashing and the upper bound}\label{sec:hashing}

The Følner sets of $\Gamma$ are doubly exponential (Proposition~\ref{prop:folner-lower}), but its chunks have far smaller models. A Følner set
records the lamp at every configuration of a large finite part of $V$. A model need only record the lamps that are lit, and it can store them
through a random affine hash of the configurations, which the affine group moves covariantly. Random $\F_2$-linear maps form a universal
hash family of Carter and Wegman~\cite{CarterWegman1979}, and a random translation does not change collision probabilities.

Let $W$ be a finite-dimensional vector space over $\F_2$, and let $K_W=S_3^W\rtimes\AGL(W)$, where $S_3^W$ is the group of functions
$W\to S_3$ under pointwise multiplication, and $g\in\AGL(W)$ acts by $(g\lambda)(w)=\lambda(g^{-1}w)$; thus $(\lambda,g)(\mu,h)=(\lambda\cdot g\mu,\,gh)$.
The support of $\lambda$ is the set of $w$ with $\lambda(w)\ne1$. Fix $k\ge1$ and a total order on $W$, let $\operatorname{Aff}(W,\F_2^k)$ be the
set of affine maps $W\to\F_2^k$, and let
\[
\Omega=\operatorname{Aff}(W,\F_2^k)\times S_3^{\F_2^k},\qquad |\Omega|=2^{k(\dim W+1)}\,6^{2^k}.
\]
For $\eta\in\operatorname{Aff}(W,\F_2^k)$ and $\lambda\in S_3^W$ let $\lambda^\eta:\F_2^k\to S_3$ send $z$ to the product, in the order of $W$, of the
$\lambda(w)$ with $\eta(w)=z$. For $\gamma=(\lambda,g)\in K_W$ define a permutation $\Theta_\gamma$ of $\Omega$ by
\[
\Theta_\gamma(\eta,u)=\bigl(\eta\circ g^{-1},\ \lambda^{\eta\circ g^{-1}}\cdot u\bigr);
\]
its inverse sends $(\eta',u')$ to $(\eta'\circ g,(\lambda^{\eta'})^{-1}u')$, and $\Theta_1$ is the identity.

\begin{lemma}[Affine hashing]\label{lem:hash}
Let $b\ge1$, let $\gamma_1=(\lambda,g)$ and $\gamma_2=(\mu,h)$ be elements of $K_W$ whose lamp parts $\lambda,\mu$ have supports of size at most
$b$, and let $\omega$ be uniformly distributed on $\Omega$. Then
\begin{enumerate}[label=(\roman*),itemsep=1pt]
\item $\Theta_{\gamma_1}\Theta_{\gamma_2}\omega=\Theta_{\gamma_1\gamma_2}\omega$ except with probability at most $\binom{2b}2\,2^{-k}$;
\item if $\gamma_1\ne\gamma_2$, then $\Theta_{\gamma_1}\omega\ne\Theta_{\gamma_2}\omega$ except with probability at most $\binom{2b}2\,2^{-k}$.
\end{enumerate}
\end{lemma}

\begin{proof}
Write $\omega=(\eta,u)$ and $\eta=\eta_0+z_0$ with $\eta_0$ linear; then $\eta_0$ and $z_0$ are independent and uniform. For $w\ne w'$ in $W$,
$\eta(w)-\eta(w')=\eta_0(w-w')$ is uniform on $\F_2^k$, because $\eta_0\mapsto\eta_0(w-w')$ is a surjective linear map; so
$\eta(w)=\eta(w')$ with probability exactly $2^{-k}$. Precomposition with a fixed affine bijection of $W$ permutes $\operatorname{Aff}(W,\F_2^k)$, so
$\eta\circ f$ is again uniform for every $f\in\AGL(W)$. Hence, for a set $S\subset W$ of size at most $2b$, $\eta\circ f$ is injective on $S$ except
with probability at most $\binom{2b}2\,2^{-k}$. If $\eta'$ is injective on a set $S$ containing the support of $\lambda$, then
$\lambda^{\eta'}(\eta'(w))=\lambda(w)$ for $w\in S$, and $\lambda^{\eta'}=1$ off $\eta'(S)$.

(i) Put $\eta'=\eta\circ(gh)^{-1}$. Then $\Theta_{\gamma_1}\Theta_{\gamma_2}(\eta,u)=(\eta',\lambda^{\eta'}\mu^{\eta\circ h^{-1}}u)$ and
$\Theta_{\gamma_1\gamma_2}(\eta,u)=(\eta',(\lambda\cdot g\mu)^{\eta'}u)$. Let $S$ be the union of the support of $\lambda$ and the image under $g$ of
the support of $\mu$, and suppose that $\eta'$ is injective on $S$. Since $\eta\circ h^{-1}(v)=\eta'(gv)$, the map $\eta\circ h^{-1}$ is injective on
the support of $\mu$, and $\mu^{\eta\circ h^{-1}}(\eta'(w))=\mu(g^{-1}w)$ for $w\in S$. The support of $\lambda\cdot g\mu$ lies in $S$. So all three
functions are trivial off $\eta'(S)$, and at $\eta'(w)$, $w\in S$, both sides equal $\lambda(w)\mu(g^{-1}w)$.

(ii) If $g\ne h$, choose $w_0$ with $h^{-1}g(w_0)\ne w_0$. If $\eta\circ g^{-1}=\eta\circ h^{-1}$, then $\eta=\eta\circ h^{-1}g$, so
$\eta(h^{-1}g(w_0))=\eta(w_0)$, which has probability $2^{-k}$. If $g=h$, then $\lambda\ne\mu$; put $\eta'=\eta\circ g^{-1}$ and let $S$ be the union
of the supports. If $\eta'$ is injective on $S$, then $\lambda^{\eta'}$ and $\mu^{\eta'}$ take the different values $\lambda(w)$ and $\mu(w)$ at
$\eta'(w)$ for some $w\in S$.
\end{proof}

To use the lemma in $\Gamma$ we need finite subgroups that contain the relevant cocycles. For $z\in\Z^2$ put $g_z=p^{z_1}q^{z_2}$, so that
$\pi(g_z)=z$, and for $m\ge1$ let $B_m=\{-m,\dots,m\}^2$. For $R\ge1$ let $Y_R\subset Y$ be the set of points of height at most $R$, let
$W_R=\F_2^{Y_R}\subset V$, and let $K_R\le\ker\pi$ consist of the elements whose lamps lie at configurations in $W_R$ and whose affine part lies in
$W_R\rtimes\GL(W_R)$, with $\GL(W_R)$ extended by the identity on the other basis vectors. Restricting to $W_R$ identifies $K_R$ with $K_{W_R}$.

\begin{lemma}[Cocycles]\label{lem:clock}
Let $m\ge1$, $R=2m+3$, $z\in B_m$ and $s\in\{p,q,t,c,a\}^{\pm1}$. Then $g_{z+\pi(s)}^{-1}sg_z$ and $g_zsg_{z+\pi(s)}^{-1}$ lie in $K_R$, and their
lamps sit at no more than one configuration.
\end{lemma}

\begin{proof}
Each of $p^{\pm1},q^{\pm1}$ changes heights by at most one and acts on the points of height at least $2$ as a translation along the rays, and these
translations commute. Hence for $(y_1,y_2)\in\Z^2$ the commutator $[p^{y_1},q^{y_2}]$ fixes every point of height greater than $|y_1|+|y_2|+1$:
while it is applied, such a point stays at height at least $2$. For $s=p^{\pm1}$ we have $\pi(s)=(\pm1,0)$, $g_{z+\pi(s)}^{-1}sg_z=1$ and
$g_zsg_{z+\pi(s)}^{-1}=p^{z_1}[q^{z_2},p^{\pm1}]p^{-z_1}$, which is supported in heights at most $|z_2|+2+|z_1|\le2m+2$. For $s=q^{\pm1}$ we have
$\pi(s)=(0,\pm1)$, $g_zsg_{z+\pi(s)}^{-1}=1$ and $g_{z+\pi(s)}^{-1}sg_z=q^{-z_2\mp1}[p^{-z_1},q^{\pm1}]q^{z_2\pm1}$, supported in heights at most
$|z_1|+2+|z_2|+1\le2m+3$. These are permutations of
$Y_R$, hence elements of $\GL(W_R)$. For $s\in\{t,c,a\}^{\pm1}$ we have $\pi(s)=0$, and $g_z^{\pm1}$ moves the points $1$ and $2$ to height at most
$2m+2$; so the conjugates of $t$ and $d$ by $g_z^{\pm1}$ are a translation and a transvection of $W_R$, while those of $a$ and $ab$ are $a$ and
$ab$, at the configuration~$0$.
\end{proof}

\begin{proposition}[Hash models]\label{prop:upper}
Let $E'$ be a chunk of $\Gamma$ each of whose elements is a product of at most $\ell\ge1$ of the generators $p,q,t,c,a$ and their inverses. Then
\[
\log\sigma_{E'}(r)\ \le\ 500\,\ell^2\,r\log(2+r)\qquad(r>1).
\]
\end{proposition}

\begin{proof}
Let $m=\lceil8\ell r\rceil$, $R=2m+3$, $D=|Y_R|=6m+9$ and $k=\lceil\log(4\ell^2r)\rceil$, and let $\Omega$ be the set above for $W=W_R$. For $g\in\Gamma$
and $z\in\Z^2$ put $\psi(g,z)=g_{z+\pi(g)}^{-1}gg_z\in\ker\pi$; then $\psi(gh,z)=\psi(g,z+\pi(h))\,\psi(h,z)$. If $g$ is a product of at most $\ell$
generators and inverses and $z\in B_{m-\ell}$, then $\psi(g,z)$ is the product of the $\psi(s,z')$ for the letters $s$ of the word, at points
$z'\in B_m$, so by Lemma~\ref{lem:clock} it lies in $K_R$ with lamps at no more than $\ell$ configurations.

Let $X=B_m\times\Omega$. For $g\in E'$ let $F_g$ be a permutation of $X$ that extends the injective partial map
$(z,\omega)\mapsto(z+\pi(g),\Theta_{\psi(g,z)}\omega)$, defined for $z\in B_{m-\ell}$, with $F_1=\mathrm{id}$. Let $g,h,gh\in E'$ and
$z\in B_{m-2\ell}$. Then $z$ and $z+\pi(h)$ lie in $B_{m-\ell}$, so $F_gF_h(z,\omega)=(z+\pi(gh),\Theta_{\psi(g,z+\pi(h))}\Theta_{\psi(h,z)}\omega)$ and
$F_{gh}(z,\omega)=(z+\pi(gh),\Theta_{\psi(g,z+\pi(h))\psi(h,z)}\omega)$; by Lemma~\ref{lem:hash}(i) these agree for all but a fraction
$\binom{2\ell}2\,2^{-k}$ of the $\omega$. Since $|B_{m-2\ell}|\ge(1-8\ell/(2m+1))|B_m|$,
\[
d\bigl(F_gF_h,F_{gh}\bigr)\ \le\ \varepsilon:=\frac{8\ell}{2m+1}+\binom{2\ell}2\,2^{-k}<\frac1{2r}+\frac1{2r}=\frac1r .
\]
If $g\ne h$ in $E'$ and $z\in B_{m-\ell}$, then either $\pi(g)\ne\pi(h)$ and $F_g(z,\omega)\ne F_h(z,\omega)$, or $\psi(g,z)\ne\psi(h,z)$ and
Lemma~\ref{lem:hash}(ii) applies; so $d(F_g,F_h)\ge1-\varepsilon$. Hence $\sigma_{E'}(r)\le|X|$.

Finally $\log|X|=2\log(2m+1)+k(D+1)+2^k\log6$. Here $2^k<8\ell^2r$, $k<\log(8\ell^2r)\le5\ell\log(2+r)$ and $D+1\le64\ell r$, so
$k(D+1)\le320\ell^2r\log(2+r)$; also $2^k\log6\le21\ell^2r\le14\ell^2r\log(2+r)$ and $2\log(2m+1)\le10\ell^2r\log(2+r)$. The sum is less than
$500\ell^2r\log(2+r)$.
\end{proof}

\subsection{The profile}

\begin{theorem}\label{thm:main-group}
Let $E\subset\Gamma$ be the chunk consisting of $1$, the five generators $p,q,t,c,a$ and their inverses, the elements represented by the prefixes
of the $25$ freely reduced words of Table~\ref{tab:certificate}, and the inverses of these elements. Let $p_H=\log107+4$ and
\[
\kappa=2+p_H=6+\log107<12.75 .
\]
There are absolute constants $c,C>0$ such that for every $r\ge2$:
\begin{enumerate}[label=(\alph*),itemsep=1pt]
\item $\log\sigma_E(r)\ge c\,r/(\log r)^\kappa$;
\item $\log\sigma_E(r)\le C\,r\log r$; more generally, every chunk $E'$ of $\Gamma$ satisfies $\log\sigma_{E'}(r)\le C_{E'}\,r\log r$.
\end{enumerate}
In particular $\sigma_E(r)=\exp(r^{1+o(1)})$. Moreover:
\begin{enumerate}[label=(\alph*),itemsep=1pt,resume]
\item For every $0<\Delta\le1$ there are $c_\Delta,e_\Delta>0$ such that the following holds for $0<e\le e_\Delta$. Let $\phi$ assign unitaries in
$U(D)$ to the five generators so that every word of Table~\ref{tab:certificate} is mapped within $e$ of $I$ in the normalized Hilbert--Schmidt
norm, and put $\phi(b)=\phi(a)\phi(c)^4$. If $\|\phi(a)\phi(b)-I\|_2\ge\Delta$, then
\[
\log D\ \ge\ c_\Delta\,\frac{e^{-2}}{(\log(1/e))^{2\kappa}},\qquad 2\kappa<25.49 .
\]
\item $\Gamma$ is finitely presented, elementary amenable and not LEF.
\end{enumerate}
\end{theorem}

\begin{proof}
\emph{Setting up the criterion.} Theorem~\ref{thm:criterion} asks for $a$ and $b$ among the generators, while in $R$ the letter $b$ abbreviates
$ac^4$. So let $S^+$ consist of the five generators and $b=ac^4\in\Gamma$, and let $R^+$ consist of $R$ and the three words $b^{-1}ac^4$, $b^2$
and $(ab)^3$ in the letters of $S^+$. These words represent $1$, $a^2,b^2,(ab)^3\in R^+$, $ab\ne1$, and the longest word of $R^+$ has length
$l=24$. The chunk $E^+$ that Theorem~\ref{thm:criterion} associates with $S^+$ and $R^+$ is contained in $E$: besides the elements of $E$, it
contains $b^{\pm1}$, the prefixes $b^{-1},c^{-4},c^{-3},c^{-2},c^{-1}$ of $b^{-1}ac^4$, and the prefixes $b$, $ab=c^4$, $aba=c^4a$, $abab=c^2$ and
$ababa=c^2a$ of $b^2$, $(ab)^3$ and $ab$. These lie in $E$, because $ac^4$, $a^2c^4$, $a^2c^4a$ and $a^2c^4a^2c^4a$ are prefixes of $b^2$ and
$(ab)^3$ in Table~\ref{tab:certificate}, $c^2$ and $c^3$ are prefixes of the relations $[d,k]$ of the family \textsf{S}, $a^2=1$, and $c^6=1$.
Restricting a model of $E$ to $E^+$ gives a model of $E^+$, so $\sigma_E\ge\sigma_{E^+}$.

For $N\ge1$ take the $M=2^N$ words $T_u$ of Theorem~\ref{thm:pair-area}, one for each configuration $u$ supported in the first $N$ points of
ray~$1$. A commutator $[T_uxT_u^{-1},T_wyT_w^{-1}]$ with $x,y\in\{a,b\}$ contains the letter $b^{\pm1}$ at most four times, and replacing each
occurrence by $(ac^4)^{\pm1}$ costs one relation $b^{-1}ac^4$. So its area with respect to $R^+$ is at most $\Lambda^+_N=\Lambda_N+4$, and
Theorem~\ref{thm:criterion} gives
\begin{equation}\label{eq:criterion-applied}
\log\sigma_E(r)\ge2^N/16\qquad\text{whenever}\qquad r\ge4.7\cdot10^{9}\cdot2^N\Lambda^+_N .
\end{equation}
Since $a\ne1$, a $(1-1/r)$-expansive map needs at least two points, so $\log\sigma_E(r)\ge1$; for $r$ below any fixed bound $r_0$ claim (a)
therefore holds once $c$ is small enough, and we may assume $r$ large.

(a) By Theorem~\ref{thm:pair-area}, $\Lambda^+_N\le C_1N^\kappa$ with $C_1$ absolute. Let $r'=r/(4.7\cdot10^9C_1)$ and
$N=\lfloor\log r'-\kappa\log\log r'\rfloor$, which is at least $1$ for large $r$. Then
$2^N\Lambda^+_N\le C_1r'(\log r')^{-\kappa}(\log r')^\kappa=r/(4.7\cdot10^9)$, so~\eqref{eq:criterion-applied} applies, and
$\log\sigma_E(r)\ge2^N/16\ge r'/(32(\log r')^\kappa)$. Since $\log r'\le\log r$, this is~(a).

(b) Every element of $E$ is a generator, the inverse of one, or represented by a prefix of a word of length at most $24$ or by its inverse.
Proposition~\ref{prop:upper} with $\ell=24$ gives $\log\sigma_E(r)\le500\cdot24^2r\log(2+r)$, and $\log(2+r)\le2\log r$ for $r\ge2$. The same
proposition applies to every chunk.

(c) Extend $\phi$ to $S^+$ by $\phi(b)=\phi(a)\phi(c)^4$. Then $b^{-1}ac^4$ is mapped to $I$ exactly, and $b^2$ and $(ab)^3$ are mapped to the
images of the words $b^2$ and $(ab)^3$ of Table~\ref{tab:certificate}; so every word of $R^+$ is mapped within $e$ of $I$. With
$\Lambda^+_N\le C_1N^\kappa$ as in (a), the condition $K_\Delta\sqrt M\,\Lambda^+_Ne\le1$ of Theorem~\ref{thm:criterion} holds when
$2^NN^{2\kappa}\le Y:=(C_1K_\Delta e)^{-2}$, which is the case for $N=\lfloor\log Y-2\kappa\log\log Y\rfloor$; this $N$ is at least $1$ for $e$ small
in terms of $\Delta$. Then $2^N\ge Y/(2(\log Y)^{2\kappa})$, and since $C_1K_\Delta\ge1$ we have $\log Y\le2\log(1/e)$. Hence
\[
\log D\ \ge\ \frac{\Delta^2}{12}\,2^N\ \ge\ \frac{\Delta^2}{24\cdot4^\kappa C_1^2K_\Delta^2}\cdot\frac{e^{-2}}{(\log(1/e))^{2\kappa}}.
\]

(d) Proposition~\ref{prop:fp} gives finite presentation and Lemma~\ref{lem:structure} elementary amenability. By (a) the profile of $E$ is unbounded, so $E$ does not embed in a finite group and
$\Gamma$ is not LEF (\cite[Lemma~2.1]{DouglasMghirbiA}). This also follows because $\Gamma$ contains $H_3$, which is finitely presented and not
residually finite, since the finitary alternating group of an infinite set is simple; a finitely presented LEF group is residually finite.
\end{proof}

\begin{remark}\label{rem:lee}
Part~(a) uses the polynomial Dehn bound of Proposition~\ref{prop:dehn}(b). With Lee's exponential bound alone, Theorem~\ref{thm:pair-area} gives
$\Lambda^+_N\le2^{BN}$ for a constant $B$ that depends on the constant in Lee's inequality, and $N=\lfloor\log(r/(4.7\cdot10^9))/(1+B)\rfloor$
in~\eqref{eq:criterion-applied} gives only $\log\sigma_E(r)\ge c\,r^{1/(1+B)}$: a superpolynomial profile, with an exponent we cannot make
explicit. Transport at scale $N$ is too expensive for an exponential Dehn bound. The grid host of Remark~\ref{rem:grid} transports at scale
$\sqrt N$ and would reach, by the argument sketched there, $\log\sigma\ge c\,r\,e^{-C\sqrt{\log r}}$ with Lee's bound alone.
\end{remark}

The hash models are far smaller than any Følner set. For $r\ge1$ let $\textup{Føl}(r)$ be the least size of a finite nonempty set
$F\subset\Gamma$ with $|Fs\setminus F|\le|F|/r$ for $s\in\{p,q,t,c,a\}$.

\begin{proposition}\label{prop:folner-lower}
For $r\ge12$,
\[
2^{\,2^{\lfloor r/12\rfloor-1}}\ \le\ \textup{Føl}(r)\ \le\ 2^{\,2^{3r+14}} .
\]
So the Følner function of $\Gamma$ is doubly exponential, while by Theorem~\ref{thm:main-group}(b) the sofic profile of every chunk is
$\exp(O(r\log r))$.
\end{proposition}

\begin{proof}
\emph{Lower bound.} Let $F$ be such a set and $l=\lfloor r/12\rfloor\ge1$. For $v\in\{0,1\}^l$ the word
$T_v=t^{v_1}p\,t^{v_2}p\cdots p\,t^{v_l}p^{1-l}$ is freely equal to $\prod_{j=1}^lp^{j-1}t^{v_j}p^{1-j}$, so it represents the translation by
$\sum_jv_je_j$, and it has length less than $3l$. The elements $a_v=T_vaT_v^{-1}$ are the copies of $a$ at $2^l$ distinct configurations, so they
are commuting involutions and form a basis of an elementary abelian group $Q$ of order $2^{M}$, $M=2^l$. Each $a_v$ is a product of fewer than
$6l$ generators and inverses; since $|Fxy\setminus F|\le|Fx\setminus F|+|Fy\setminus F|$ and $|Fx^{-1}\setminus F|=|Fx\setminus F|$, we get
$|Fa_v\setminus F|\le6l|F|/r\le|F|/2$, and so $\sum_v|Fa_v\setminus F|\le M|F|/2$.

On the other hand, for every finite $S\subset Q$,
\begin{equation}\label{eq:cube}
\sum_v|Sa_v\setminus S|\ \ge\ |S|\,(M-\log|S|).
\end{equation}
Indeed, the left side is $M|S|-2e(S)$, where $e(S)$ is the number of pairs $\{x,xa_v\}\subset S$, and $e(S)\le\frac12|S|\log|S|$ by the
edge-isoperimetric inequality of the cube~\cite{Harper1964}. We recall the proof, by induction on $M$: splitting $S$ by the last coordinate into
$S_0,S_1$ of sizes $s_0\le s_1$ gives $e(S)\le e(S_0)+e(S_1)+s_0\le\frac12(s_0\log s_0+s_1\log s_1)+s_0$, which is at most $\frac12s\log s$,
$s=s_0+s_1$, because the binary entropy function $h$ satisfies $h(x)\ge2x$ for $0\le x\le\frac12$. The right multiplications by the $a_v$ preserve each coset $xQ$;
applying~\eqref{eq:cube} to the sets $x^{-1}(F\cap xQ)$ and using $|F\cap xQ|\le|F|$ gives $\sum_v|Fa_v\setminus F|\ge|F|(M-\log|F|)$. Hence
$\log|F|\ge M/2=2^{l-1}$.

\emph{Upper bound.} Let $m=\lceil(r-1)/2\rceil$, $R=2m+3$ and $F=\bigcup_{z\in B_m}K_Rg_z$, a disjoint union because $\pi(g_z)=z$. By
Lemma~\ref{lem:clock}, $g_zs\in K_Rg_{z+\pi(s)}$ for $z\in B_m$ and each generator $s$, so $|Fs\setminus F|\le|K_R|(2m+1)=|F|/(2m+1)\le|F|/r$. With
$D=|Y_R|=6m+9$ we have $|K_R|=6^{2^D}\cdot2^D\cdot|\GL_D(\F_2)|$, so $\log|F|\le2^D\log6+D+D^2+2\log(2m+1)\le2^{D+2}\le2^{3r+14}$.
\end{proof}

\subsection{The grid host}\label{sec:grid}

\begin{remark}\label{rem:grid}
Placing the lamps on configurations of the cells of a grid, rather than of the points of $Y$, should give a stronger lower bound with the same upper
bound, and a nearly linear exponent $r^{1-o(1)}$ even from Lee's exponential Dehn bound alone. We only sketch the proofs, and the conclusions of this remark are conditional on them. Let $I=Y\times Y$, let $T=H\times H$ act on $I$ coordinatewise, and let
\[
\Gamma_2=\Bigl(\bigoplus_{v\in V_2}S_3\Bigr)\rtimes\bigl(V_2\rtimes(\GL_{\rm fin}(I)\rtimes T)\bigr),\qquad V_2=\F_2^{(I)} .
\]
It is locally finite-by-$\Z^4$ and generated by seven elements (commutators $[E_{x\leftarrow y},E_{y\leftarrow z}]=E_{x\leftarrow z}$ of
transvections of the third kind give those within a row or a column): the generators $p_1,q_1,p_2,q_2$ of the two copies of $H$, the translation $t$ by
the cell $o=((1,1),(1,1))$, the element $c=d\,ab$, where $d$ is the transvection that adds $e_o$ to $e_{((1,2),(1,2))}$, and $a$. A certificate of
$77$ relations, of total length $5445$ and maximum length $224$, is listed in \texttt{verification/relators.txt} and checked by
\texttt{verification/relators\_check.py}. It has the families of Table~\ref{tab:certificate}, with two copies of the Houghton relators and
their commutation, and with stabilizer and action relations for three kinds of transvection: within a row, within a column, and neither.

The gain is in the scale of transport. $N$ cells fit in a square of side $\lceil\sqrt N\,\rceil$, so the transporters have length $O(\sqrt N)$, and
the local fillings cost $C(\delta_H(C\sqrt N)+N)$; the term $N$ counts the swaps of letters from the two copies of $H$ in a change of transporter.
The clearing argument of Theorem~\ref{thm:pair-area} then bounds the pair areas by $CN^2(\delta_H(C\sqrt N)+N)$, which is at most $CN^{2+p_H/2}$ by
Proposition~\ref{prop:dehn}(b) and at most $e^{C\sqrt N}$ by Proposition~\ref{prop:dehn}(a). As in the proof of Theorem~\ref{thm:main-group}, the chunk
$E_2$ of this certificate should therefore have
\[
\log\sigma_{E_2}(r)\ \ge\ c\,\frac{r}{(\log r)^{2+p_H/2}},\qquad 2+\tfrac{p_H}2<7.38,
\]
and, using only Lee's bound, $\log\sigma_{E_2}(r)\ge c\,r\,e^{-C\sqrt{\log r}}$; unitary models need $\log D\ge c_\Delta e^{-2}/(\log(1/e))^{4+p_H}$.
In the other direction, hashing the configurations directly, as in Proposition~\ref{prop:upper}, would need a configuration space of dimension
of order $r^2$. Hashing the cells first avoids this: a cell $(y_1,y_2)$ gets the label $a_1(f_1^{-1}y_1)+a_2(f_2^{-1}y_2)\in\F_q$, where the $f_j$ are
positions in finite F{\o}lner charts of $H$ of size $e^{O(r\log r)}$ (inverses of right F{\o}lner sets, since the charts move by left
multiplication) and the $a_j$ are random tables on $O(r)$ points; this label is covariant under the
charts and separates the finitely many cells of a chunk with probability $1-O(1/q)$. The configurations are then compressed linearly to
$\F_2^{\F_q}$ and hashed affinely as before. With $q$ of order $\ell^2r$, for a chunk of words of length at most $\ell$, and the other parameters of order $r$, this would give $\log\sigma_{E'}(r)\le C_{E'}r\log(2+r)$ for
every chunk $E'$ of $\Gamma_2$, and the same for the $d$-dimensional grid host, built in the same way on the cells of $Y^d$ with $H^d$ acting,
for every fixed $d$. So $\sigma_{E_2}(r)$ would be $\exp(r^{1+o(1)})$ as well, and no fixed-dimensional
grid host would exceed $\exp(r^{1+o(1)})$. The proofs follow those above with the changes indicated, and we omit them.
\end{remark}
